\documentclass[11pt,a4paper]{report}

%INPUT of SETTINGS (usepackage, styles, new command,...)
\usepackage{geometry}
\geometry{tmargin=2cm,bmargin=2cm,lmargin=3cm,rmargin=3cm,includefoot, includehead}
\usepackage[english]{babel}
\usepackage[T1]{fontenc}
\usepackage[latin1]{inputenc}
\usepackage{graphicx}
\usepackage{ blindtext, color}
\usepackage{fancyhdr}
\usepackage{hyperref}
\usepackage{float}
\usepackage{enumitem}
\usepackage{lastpage}
\usepackage{listings}
\usepackage{color}
\usepackage{amsmath}
\usepackage{tabularx}
\usepackage{array}
\usepackage{verbatim}
\usepackage{booktabs}
\usepackage{ltxtable} 
\usepackage{multirow}
\usepackage{titlesec}
\makeatletter\let\MPtrue\@minipagetrue\makeatother

%----------------------------------------------------------------------------------------
%   Title Style Chapter
%----------------------------------------------------------------------------------------
\definecolor{gray75}{gray}{0.70}
\newcommand{\hsp}{\hspace{15pt}}
\newcommand\tab[1][1cm]{\hspace*{#1}}
\newcommand{\forceindent}{\leavevmode{\parindent=1em\indent}}
\titleformat{\chapter}[hang]{\Huge\bfseries}{\thechapter\hsp\textcolor{gray75}{|}\hsp}{0pt}{\Huge\bfseries}


%----------------------------------------------------------------------------------------
%  ToC
%----------------------------------------------------------------------------------------
\hypersetup{
    colorlinks,
    citecolor=black,
    filecolor=black,
    linkcolor=black,
    urlcolor=black
}

% Renaming of ToC
\addto\captionsenglish{% Replace "english" with the language you use
  \renewcommand{\contentsname}%
    {Table of Contents}%
}

%----------------------------------------------------------------------------------------
%   Header/Footer
%----------------------------------------------------------------------------------------
\pagestyle{fancy} 
\fancyhf{}
%\renewcommand{\chaptermark}[1]{\markboth{#1}{}}

%\renewcommand{\chaptermark}[1]{%
% \markboth{\chaptername
% \ \thechapter.\ #1}{}}

%\renewcommand{\sectionmark}[1]{\markright{\thesection.\ #1}}

\lhead{\nouppercase{\leftmark}}
%\chead{}
\rhead{\nouppercase{\rightmark}}
\lfoot{PPD - \textit{Version 1.0}}
%\cfoot{}
\rfoot{\thepage}
\renewcommand{\headrulewidth}{0.008pt}
\renewcommand{\footrulewidth}{0.008pt}

% Redefine the plain page style
\fancypagestyle{plain}{%
	\fancyhf{}%
	\fancyfoot[R]{\thepage}%
	\renewcommand{\headrulewidth}{0pt}% Line at the header invisible
	\renewcommand{\footrulewidth}{0.008pt}% Line at the footer visible
}

%-----------
% New Page for each section
%---------
%\newcommand{\sectionbreak}{\clearpage}

%-------
% Numbered list for goals
%-------
\newlist{Goal}{enumerate}{1}
\setlist[Goal]{label=[G\arabic*],leftmargin=0.6in}

%-------
% Numbered list for requirements
%-------
\newlist{Req}{enumerate}{1}
\setlist[Req]{label=[R\arabic*],leftmargin=0.6in}

%-------
% Numbered list for domain assumptions
%-------
\newlist{DA}{enumerate}{1}
\setlist[DA]{label=[DA\arabic*],leftmargin=0.6in}

%-------
% Numbered list for text assumptions
%-------
\newlist{TA}{enumerate}{1}
\setlist[TA]{label=[TA\arabic*],leftmargin=0.6in}

%-----
% Image box, arguments: file path, caption
%----
\newcommand{\imagebox}[4]{
\begin{figure}[H]
\centering
\fbox{\includegraphics[width=\textwidth,keepaspectratio,angle=#3, scale=#4]{#1}}
\caption{#2}
\end{figure}
}

\newcommand{\imageboxNoFit}[4]{
\begin{figure}[H]
\centering
\fbox{\includegraphics[keepaspectratio,angle=#3,scale=#4]{#1}}
\caption{#2}
\end{figure}
}

%---
% Alloy listing
%---
\definecolor{pgreen}{rgb}{0,0.64,0.013}
\definecolor{pblue}{rgb}{0.13,0.13,1}

\lstdefinelanguage{pseudocode}{
	keywords={
		function, if, endif, for, each, endfor, return, called, in, boolean, integer, float, Position, Car, List, PositionMediaVariance, not, break},
	literate=%
	{:}{$\colon$}1
	%{|}{$\bullet$}1
	%{==}{$=$}1
	%{=}{$=$}1
	{!=}{$\neq$}1
	%{&&}{$\land$}1
	%{||}{$\lor$}1
	{<=}{$\le$}1
	{>=}{$\ge$}1
	%{all}{$\forall$}1
	%{exists}{$\exists$}1
	{!in}{$\not\in$}1
	{\\in}{$\in$}1
	{=>}{$\implies$}3
	% the following isn't actually Alloy, but it gives the option to produce nicer latex
	{|=>}{$\Rightarrow$}2
	{<=set}{$\subseteq$}1
	{+set}{$\cup$}1
	{*set}{$\cap$}1
	{==>}{$\Longrightarrow$}3
	{<=>}{$\iff$}4
	{...}{$\ldots$}1
	{\\hl}{$\hline$}1
	{\\alpha}{$\alpha$}1
	{\\beta}{$\beta$}1
	{\\gamma}{$\gamma$}1
	{\\delta}{$\delta$}1
	{\\epsilon}{$\epsilon$}1
	{\\zeta}{$\zeta$}1
	{\\eta}{$\eta$}1
	{\\theta}{$\theta$}1
	{\\iota}{$\iota$}1
	{\\kappa}{$\kappa$}1
	{\\lambda}{$\lambda$}1
	{\\mu}{$\mu$}1
	{\\nu}{$\nu$}1
	{\\xi}{$\xi$}1
	{\\pi}{$\pi$}1
	{\\rho}{$\rho$}1
	{\\sigma}{$\sigma$}1
	{\\tau}{$\tau$}1
	{\\upsilon}{$\upsilon$}1
	{\\phi}{$\phi$}1
	{\\chi}{$\chi$}1
	{\\psi}{$\psi$}1
	{\\omega}{$\omega$}1
	{\\Gamma}{$\Gamma$}1
	{\\Delta}{$\Delta$}1
	{\\Theta}{$\Theta$}1
	{\\Lambda}{$\Lambda$}1
	{\\Xi}{$\Xi$}1
	{\\Pi}{$\Pi$}1
	{\\Sigma}{$\Sigma$}1
	{\\Upsilon}{$\Upsilon$}1
	{\\Phi}{$\Phi$}1
	{\\Psi}{$\Psi$}1
	{\\Omega}{$\Omega$}1
	{\\EOF}{\;}1,
	sensitive=true,
	morecomment=[l]//,%
	%morecomment=[l]{--},%
	morecomment=[s]{/*}{*/},%
	morestring=[b]",
	numbers=none,
	firstnumber=1,
	%numberstyle=\tiny,
	%basicstyle=\scriptsize\ttfamily,
	commentstyle=\color{pgreen},
	keywordstyle=\color{pblue},
	%ndkeywordstyle=\bfseries,
}

\lstnewenvironment{pseudocodeEnv}{%	
	\lstset{language=pseudocode,
		floatplacement={tbp},
		captionpos=b,
		xleftmargin=2pt,
		xrightmargin=2pt,
		breaklines=true,
		tabsize=3,
		%basicstyle=\ttfamily,
	}
} {}
%--- End of listing for Alloy
%
%---
% Recall item commands
%---
\newcommand\itemtext[2]{%
	\expandafter\gdef\csname item#1\endcsname{#2}%
	\label{#1}#2}

\newcommand\useitemtext[2]{\csname item#1\endcsname}

%--- 
% Table for use case descriptions
%---
\newcommand{\usecasedesc}[6]{%
	\begin{table}[H]
		\begin{tabular}{p{0.3\linewidth} p{0.7\linewidth}}
			%\begin{tabular}{\linewidth}{>{\setlength\hsize{.25\hsize}}X>{\setlength\hsize{.75\hsize}\setlength\linewidth{\hsize}}X}
			\textit{Name} & #1 \\
			\textit{Actors} & #2 \\
			\textit{Entry conditions} &\MPtrue#3\\
			\textit{Flow of events} &\MPtrue#4\\
			\textit{Exit conditions} & #5\\
			\textit{Exceptions} & #6\\
		\end{tabular}
	\end{table}
}

%\newcommand{optusecase}{[leftmargin=0.50cm,nosep]}


%\setlength{\abovecaptionskip}{3pt}

%-------------------------------------------------------------------------
%   START DOCUMENT 
%-------------------------------------------------------------------------
\begin{document}

% INPUT OF TITLE PAGE
 %----------------------------------------------------------------------------------------
%	TITLEPAGE
%----------------------------------------------------------------------------------------
\begin{titlepage}

\newcommand{\HRule}{\rule{\linewidth}{0.3mm}} % Defines a new command for the horizontal lines, change thickness here

%\vspace*{0.3cm}
 
\center % Center everything on the page

%----------------------------------------------------------------------------------------
%	LOGO SECTION
%----------------------------------------------------------------------------------------
\includegraphics[scale=0.33]{../../logo.png}\\[0.4cm] 

%----------------------------------------------------------------------------------------
%	HEADING SECTION
%----------------------------------------------------------------------------------------
\textsc{\Large \textbf{Politecnico di Milano}\\MSc Computer Science and Engineering}\\[1.7cm] \textsc{\Large \textbf{Software Engineering 2}\\[0.1cm]Academic Year 2016-2017}\\[1cm] 
%\textsc{\LARGE Minor Heading}\\[0.5cm] 

%----------------------------------------------------------------------------------------
%	TITLE SECTION
%----------------------------------------------------------------------------------------

\HRule \\[0.4cm]
{ \huge \bfseries Project Plan Document\\[0.5cm]\textit{PowerEnJoy}}\\[0.4cm] % Title of document
\HRule \\[1.5cm]
 
%----------------------------------------------------------------------------------------
%	AUTHOR SECTION
%----------------------------------------------------------------------------------------
\Large \emph{Authors:}\\
\LARGE{
\begin{center}
\begin{tabular}{ r  l }
	Melloni Giulio & 876279\\
	Renzi Marco & 878269\\
	Testa Filippo & 875456 \\
\end{tabular}
\end{center}
}
\vspace*{0.8cm}
\Large{\emph{Reference Professor:} \\
{\textsc{Mottola} Luca} 
}

%----------------------------------------------------------------------------------------
%	DATE SECTION
%----------------------------------------------------------------------------------------
 \vspace*{1cm}
{\LARGE {\it Release Date: January 22$^{th}$, 2017\\Version 1.0}}
%vfill % Fill the rest of the page with whitespace

\end{titlepage}

%-------------------------------------------------------------------------
%   TOC
%-------------------------------------------------------------------------
%\pagestyle{plain}
\pagenumbering{roman}
\thispagestyle{empty}
\tableofcontents
\listoffigures
\cleardoublepage
\pagenumbering{arabic}
\pagestyle{fancy}

%--------------------------------------------------------------------
%  BEGIN TEXT
%--------------------------------------------------------------------
\chapter{Introduction}
\section{Revision History}
\begin{tabularx}{\textwidth}{|l|l|X|}
\hline
Version & Date & Summary\\
\hline
1.0 & 22/01/2017 & First release of the document\\
\hline
\end{tabularx}
\section{Purpose and Scope}
\section{Definitions, Acronyms, Abbreviations}
\begin{description}
	\item[PowerEnJoy]is the name of the system that has to be developed.
	\item[System]sometimes called also \textit{system-to-be}, represents the application that will be described and implemented.
	In particular, its structure and implementation will be explained in the following documents. People that will use the car-sharing service will interact with it, via some interfaces, in order to complete some operations (e.g.: reservation and renting).
	\item[Renting]it is the act of picking-up an available car and of starting to drive.
	\item[Ride] the event of picking-up a car, driving through the city and parking it. Every Ride is associated to a single user and to a single car.
	\item[Reservation]it is the action of booking an available car.
	\item[Car] a car is an electrical vehicle that will be used by a registered user.
	\item[Not Registered User] indicates a person who hasn't registered to the system yet; for this reason he can't access to any of the offered function. The only possible action that he can carry out is the registration to get a personal account.
	\item[Registered User] interacts with the system to use the sharing service. He has an account (which contains personal information, driving license number and payment data) that must be used to access to the application in order to exploit all the functionalities.
	\item[Employee] it's a person who works for the company, whose main task is to plug into the power grid those cars that haven't been plugged in by the users. He is also in charge of taking care of the status of the cars and of moving the vehicles from a safe area to a charging area and vice versa if needed.
	\item[Safe Area] indicates a set of parking lots where the users have to leave the car at the end of the rent; the set of the Safe Areas is pre-defined by the system management. These areas are spread all over the city.
	\item[Plug] defines the electrical component that physically connects the car to the power grid.
	\item[Charging Area] is a special \textit{Safe Area} that also provides a certain number of plugs that connect the cars to the power grid in order to recharge the battery.
	\item[Registration] the procedure that an unregistered user has to perform to become a registered user. At the end, the unregistered user will have an account. To complete this operation three different types of data are required: personal information, driving license number and payment info.
	\item[Search] this functionality lets the registered user search for available cars within a certain range from his/her current position or from a specified address.
	\item[RASD] is the acronym of \textit{Requirements Analysis and Specification Document}
	\item[DD] is the acronym of \textit{Design Document}
	\item[ITPD] is the acronym of \textit{Integration Test Plan Document}
	\item[PPD] is the acronym of \textit{Project Plan Document}
\end{description}
\section{Reference Documents}
	\begin{itemize}
	\item Project Assignments 2016-2017
	\item RASD v1.1
	\item DD v1.1
	\item ITPD v1.0
\end{itemize}
\chapter{Project Size Estimation}
\chapter{Project Cost and Effort Estimation}
\chapter{Schedule}
\chapter{Resource Allocation}
\chapter{Risk Management and Mitigation}
This section deals with the possible risks that can threaten the system during its development process. These hazards can stem from a variety of sources and thus can be classified as:
\begin{itemize}
	\item \textbf{Project risks}: if they threaten the project plan, such as the project schedule. As can be seen in the following tables, the most common problem that can arise from such risks is a delay in the project release.
	\item \textbf{Technical risks}: they are related to the technical part of the project. They include a major variety of problems such as unskilled staff, flaws in the external adopted components, changes in the set of the requirements and so on.
	\item \textbf{Business risks}: they include a set of heterogeneous hazards such as budget cuts, sales falls and market policy flaws. If they become concrete, they can compromise the viability of the project.
	\item \textbf{Personnel risks}: this type of hazards deals with all the possible problems that can be met within the team of work.
\end{itemize}
Since a pro-active risk-management approach is desirable, for each category of risk a list of the most recurrent hazards will be taken into account, along with their probabilities to be faced, their levels of impact on the project and the strategies that can be adopted to mitigate them.
\section{Project risks}
    \begin{tabularx}{\textwidth}{ |p{3cm}|p{1.8cm}|p{1.85cm}|X|}
    \hline
    Risk & Probability & Impact & Strategy \\
    \hline
    Underestimated development time & Moderate & Serious & If the spare time is not sufficient, try to negotiate with the customer the possibility to have two releases, the first to the planned one and the second one as early as possible. \\ \hline
    A change in the direction of the project & Low & Moderate & Redact very precise and detailed documentation of the project so that new managers are able to handle the process. \\ \hline
    \end{tabularx}
\section{Technical risks}
    \begin{tabularx}{\textwidth}{ |p{3cm}|p{1.8cm}|p{1.85cm}|X|}
    \hline
    Risk & Probability & Impact & Strategy \\ \hline
    Difficulty in recruiting a skilled staff & Moderate & Serious & Be prepared and inclined to pay extra to find skilled people. If it is not sufficient, may consider the possibility of buying external components that are already developed. \\ \hline
    Changes in the requirements & Low & Serious & Be compliant with the requirements traced in the \textit{RASD} and in case changes occur evaluate the trade off related to the needed variations in the project. Increment on the price should come from the customer. \\ \hline
    External components have flaws & Moderate & Serious & Choose components that are on the market for a long time, so that they can be reliable. Experience with precedent projects may be an advantage in this case. \\ \hline
    Modification in the APIs of the external components & Low & Moderate & Write the code as portable as possible. Changes in the APIs should not break the system. Also plan updates in case such modifications occur. \\ \hline
    \end{tabularx}
\section{Business risks}
    \begin{tabularx}{\textwidth}{ |p{3cm}|p{1.8cm}|p{1.85cm}|X|}
    \hline
    Risk & Probability & Impact & Strategy \\ \hline
    Budget cuts during the development & Moderate & Tremendous & This kind of risk is the most critical one. In order to mitigate it, redact a document in which the financial benefits associated to the business of the project is clearly stated and try to convince managers not to apply cuts. \\ \hline
    Poor sales & Moderate & Serious & Make special discounts and offers at the launch of the product in order to achieve great popularity. In doing so, a market analysis has to be carried out . \\ \hline
    Competitors in the market & Moderate & Moderate & Choose a guide line for your product: should it be cheaper than the other or should it aim to optimality? In the launch phase also consider the possibility of making special discounts to attract clients. \\ \hline
    \end{tabularx}
\section{Personnel risks}
    \begin{tabularx}{\textwidth}{ |p{3cm}|p{1.8cm}|p{1.85cm}|X|}
    \hline
    Risk & Probability & Impact & Strategy \\ \hline
    Poor motivation & Moderate & Moderate & Try to select people that are skilled in particular fields and have experience in their work. Praise their commitment with rewards. \\ \hline
    Conflicts within the work team & Low & Moderate & Promote the cooperation of people by making a community that periodically meets and discuss on the level of the project. \\ \hline
    Key staff are ill at critical times in the project & Moderate & Serious & Consider the possibility to have several people working on common tasks in order to prevent excessive time losses. \\ \hline
    \end{tabularx}
\chapter{Effort Spent}
%TODO complete hours
In order to complete this document, each author worked for XX hours.
%-------------------------------------------------------------------------
%   END DOCUMENT
%-------------------------------------------------------------------------
\end{document} 